% Pursuing Stacks (cotes 134-2 to 134-8): archive cote 134-5, site page 29 (PDF page 30), typescript page 216.
% typescript. Draft compared: ps4.tex:1359-1429.
% AI-transcribed by Claude Opus 5.5 (claude-opus-5-5, adaptive thinking, effort high), Message Batches API, batch pilot (msgbatch_0189Nhf1xB615YWnztubtLfc), 2026-11-12; original transcription was not yet checked by eye; reviewed in this edition, see page ledger.
\scanpage{134-5/29}
\mspage{216}

(vi') Same as (vi), with F restricted to be an object b in B.

Indeed, we have of course (i)$\hand{\Rightarrow}$(vi)$\hand{\Rightarrow}$(vi'), but also (vi')$\hand{\Rightarrow}$(iii)
if we admit strong saturation of $\underline{W}$, which implies that an object X of
(Cat) is $\underline{W}$-aspheric iff it's image in Hot($\underline{W}$) is a final object.

\hand{\uline{\textcolor{black}{Remark}}} 3. If we don't assume A and B to be $\underline{W}$-aspheric, the implications
\[ \text{(As)} \hand{\Longrightarrow} \text{(iv')} \hand{\Longrightarrow} \text{(iv)} \]
are both strict. As an illustration of this point, take for B a discre-%
te category (defined in terms of a set of indices I = Ob B), thus a
\margin{$(A_b)_{b\in I}$} \xout{functor of a} category A over B is essentially the same as a family $(A_{\hand{b}})_{\hand{b\in I}}$
of objects of (Cat) indexed by I. In terms of this\xout{e} family, we see at
once that the three conditions above on i: A$\hand{\rightarrow}$B mean respectively
(a) that all categories $A_b$ are $\underline{W}$-aspheric, (b) that all categories
$A_b$ are non-empty, and (c) condition vacuous. This example brings to
mind that \hand{\uline{\textcolor{black}{condition}}} (iv) \hand{\uline{\textcolor{black}{is a Serre-fibration type condition}}} , we'll
come back upon this condition when studying \xout{\illegible{1}} homotopy properties
of maps in (Cat)\hand{,} with special emphasis on base change questions. Like-%
wise, condition (iv') appears as a strengthening of such Serre-type
condition, to the effect that the restriction of A over any connected
component of B should be moreover non-empty. As an example, we may take
the projection BxC $\hand{\rightarrow}$ \over{\unsure{C}}{B} , where C is any non-empty object in (Cat).

I would like now to dwell still a little on the case, of special
interest of course for the modelizing story, when A and B are test cate-%
gories or something of the kind. The formulation (i) of the asphericity
condition \over{of}{for} the functor i can be expressed by stating that in the
the functor $\hand{\overline{\textcolor{black}{i^{*}}}}$ \xout{on} \inserted[typed]{between} the localizations $\mathrm{Hot}_B$ and $\mathrm{Hot}_A$ exists, and gives
rise (via \hand{$\lambda$}$_i$) to a commutation morphism which is an \hand{\uline{\textcolor{black}{isomorphism}}}
% Hand-drawn arrows are marked blue; typed objects remain plain.
\diagram{
\begin{tikzcd}[column sep=small]
\mathrm{Hot}_B \arrow[color=blue, rr, "\hand{\overline{\textcolor{black}{i}^{*}}}"] \arrow[color=blue, dr, swap, "\hand{\overline{i_B}}", ""{name=iB, right}] & & \mathrm{Hot}_A \arrow[color=blue, dl, "\hand{\overline{i_A}}", ""{name=iA, left}] \\
& \mathrm{Hot}(\underline{W}) \arrow[color=blue, bend right=10, from=iA, to=iB, "\hand{\overline{\lambda_i}}"', "\hand{\sim}"] &
\end{tikzcd}
}
\margin{$\overset{\overline{\lambda_i}}{\underset{\sim}{\longleftarrow}}$}
.

This shows, when A and B are pseudo-test categories for $\underline{W}$ , i.e. the
functors $\hand{\overline{\textcolor{black}{i_A}}}$ and $\hand{\overline{\textcolor{black}{i_B}}}$ are equivalences, that the functor $\hand{\overline{\textcolor{black}{i^{*}}}}$ deduced
from the \inserted{$\underline{W}$-}aspheric \xout{functor} map i: A$\hand{\rightarrow}$B , \hand{\uline{\textcolor{black}{does not depend}}} (up to cano-%
nical isomorphism) on the choice of i , and can be described as the com-%
position of $\hand{\overline{\textcolor{black}{i_B}}}$ followed by a quasi-inverse for $\hand{\overline{\textcolor{black}{i_A}}}$ . This of course
is very much in keeping with the ``inspiring assumption'' (section 28),
which just means that up to unique isomorphism, there is indeed but \hand{\uline{\textcolor{black}{one}}}

\margin{29}
